% Generated draft: requires mathematical and source review.
\documentclass[11pt]{article}
\usepackage[margin=25mm]{geometry}
\usepackage{fontspec}
\setmainfont{Latin Modern Roman}
\usepackage{amsmath,amssymb,mathtools,unicode-math}
\setmathfont{Latin Modern Math}
\usepackage{enumitem,needspace,xurl,hyperref}
\setlength{\parindent}{0pt}
\setlength{\parskip}{5pt}
\setlength{\emergencystretch}{3em}
\newcommand{\N}{\mathbb{N}}
\newcommand{\Z}{\mathbb{Z}}
\newcommand{\Q}{\mathbb{Q}}
\newcommand{\R}{\mathbb{R}}
\newcommand{\C}{\mathbb{C}}
\newcommand{\F}{\mathbb{F}}
\newcommand{\A}{\mathbb{A}}
\newcommand{\PP}{\mathbb{P}}
\newcommand{\E}{\mathbb{E}}
\newcommand{\eps}{\varepsilon}
\newcommand{\dd}{\,\mathrm{d}}
\begin{document}
\section*{UnsolvedMath ID: 30000687}
\textbf{Source problem code:} catalogue entry
\subsection{Definitions and mathematical statement}
Let $(W,S)$ be a Coxeter system and let $\theta$ be the intended involutive automorphism, with the usual Coxeter-system convention $\theta(S)=S$. Impose exactly the source's no-odd-order condition:
\[
\theta(s)\ne s\quad\Longrightarrow\quad
\operatorname{ord}(s\theta(s))\text{ is even or infinite}
\qquad(s\in S).
\]
Thus infinite order is permitted, and the order-one product for a fixed generator is not excluded.

Use the standard set of twisted identities
\[
\iota(\theta)=\{\theta(w^{-1})w:w\in W\},
\]
and let $\operatorname{Br}(\iota(\theta))$ be this set with the order induced from Bruhat order on $W$. For comparable $u<v$ in this poset, let $\Delta(u,v)$ be the order complex of the open interval: its simplices are chains $u<x_1<\cdots<x_k<v$. This is the usual topological interpretation of an interval in the stated sphere/ball question; the order complex of the whole closed interval would instead have cone points.

The conjectural strengthening asks whether every such interval complex is either a PL sphere or a ball. A PL sphere is a simplicial complex piecewise-linearly homeomorphic to a standard sphere; a ball is at least required to be homeomorphic to a closed Euclidean ball. The source does not explicitly say whether the latter must also be a PL ball.

The original statement reports Theorem 7 as giving the weaker conclusion ``PL sphere or $\Z$-acyclic'', where acyclic means vanishing reduced integral homology in every degree. The requested strengthening replaces that acyclic alternative by a ball; acyclicity alone is not identified with the ball property. The extract does not state the endpoint/empty-complex dimension conventions or explicitly define the twisted-identity notation and preservation of $S$, so the standard interpretations above remain subject to source confirmation.
\subsection{Short English statement}
Under the specified no-odd-order condition, must every twisted-identity Bruhat interval have the topology of a PL sphere or a ball? The source reports a sphere-or-integrally-acyclic theorem and asks to strengthen the acyclic alternative to a ball.
\subsection{Sources}
\begin{itemize}
\item Supplied report: \emph{Geometric and Topological Combinatorics} (2007), Oberwolfach Reports, DOI 10.4171/owr/2007/04, pp. 195--272; proposed year 2007. \url{https://doi.org/10.4171/owr/2007/04}.
\item Local UnsolvedMath registry, permanent ID 30000687, catalogue entry; created 25 August 2026, updated 16 September 2026. The original record names Theorem 7 and its PL-sphere-or-$\Z$-acyclic conclusion. The corrected\_verified note says the proposed ball strengthening and hypothesis were retained while the theorem-number reference was removed from the clean statement. The supplied 21 August 2026 literature triage labels the strengthening open and reports no located authoritative complete resolution. No additional external reference is supplied; no independent literature assessment is asserted here.
\end{itemize}
\end{document}
